% Replacement for the Psalm 64 paragraph in the combined Gradual/Alleluia
% subsection of main.tex. Kept separate while the component revisions are
% integrated serially.
The Alleluia takes only the opening of Psalm 64, but that opening announces a
larger movement: praise and a fulfilled vow at Zion lead into heard prayer,
pardon for overpowering sins, the good things of God's house, and finally a
creation so abundantly watered that valleys and fields sing. The appointed
verse therefore names neither private admiration nor a merely interior
resolution. Its praise belongs to assembled worship, and its vow is something
rendered to God. The inherited Greek and Latin title associates the psalm with
return from exile; that title is a traditional interpretive horizon, not proof
of authorship, date, or the historical occasion of the Hebrew psalm.

Augustine and Bellarmine preserve different stages of a literal-to-spiritual
movement. Augustine reads Zion and Jerusalem as the heavenly homeland toward
which Christian hope is already anchored. He develops the vow
eschatologically: what is begun now is rendered whole in the resurrection,
soul and incorruptible body together, as a total offering kindled by charity
(\emph{Enarratio in Psalmum 64} 3--4). Bellarmine begins instead with the
people's desire to return from Babylon to the Temple, where praise and
sacrificial vows are rightly offered, and only then argues \emph{a fortiori}
for the heavenly Jerusalem, whose citizens offer themselves in ardent love
(\emph{Commentary on Psalm 64}, explanation 1--2). These are Christian
receptions of the psalm, not replacements for its canonical horizon.
Augustine's anti-Jewish and supersessionist rhetoric in the immediately
preceding sections is neither repeated nor made to support the retained
exegesis.

Cassiodorus supplies an unusually close reception of the psalm's full context.
Among the ``good things'' of God's house he names the diverse gifts of
1~Corinthians~12---wisdom, knowledge, faith, healings, tongues, and
interpretations. When the psalm's valleys later abound with grain, he names
Luke's publican: the man lowers his eyes, refuses presumption in merits,
confesses his weakness, and is raised by the Lord after humbling himself
(\emph{Expositio Psalmorum}, on Psalm 64, local verses 5 and 14). Cassiodorus
thus lets the larger psalm illuminate both the Epistle and the Gospel. That is
a documented later correspondence, not evidence that a Roman compiler chose
or arranged these chants for this reason. Within the appointed sequence, the
Gradual's appeal to God's equitable sight and the Alleluia's God-directed
praise move away from self-certifying judgment; that final observation remains
editorial synthesis rather than recovered intent.
